\begin{textsample}{POS Dim 1 – pa – Score 70.00 – pa/pa\_em\_20.txt}  \label{ex:f1_pos_032}
b ) Campo de Saberes e Práticas do Ensino da Arte O Campo de Saberes e Práticas do Ensino da Arte integra -se à Área de Linguagens e suas Tecnologias, por meio de suas quatro linguagens : Artes Visuais, Música, Teatro e Dança.
Essas passam a se articular para a produção de saberes \textbf{que} envolvem a compreensão do fazer artístico, posto \textbf{que} a Arte é \textbf{uma} área de conhecimento \textbf{que} \textbf{opera}, segundo Aristóteles ( 1987, Passim ), por meio de \textbf{uma} “ poiesis ”, isto é, \textbf{um} ato, \textbf{uma} ação criativa na qual as \textbf{fronteiras} entre quem faz ( sujeito/criador ) e o \textbf{que} faz ( objeto/criação ) podem variar de acordo com a linguagem artística em questão, chegando mesmo a linhas muito tênues quando se trata da linguagem do Teatro e das Artes Performáticas \textbf{Contemporâneas}, momento em \textbf{que} é inviável discernir o “ sujeito/criador ” do “ objeto/criação ”.
Todo esse processo está envolto na construção do conhecimento em Arte, \textbf{que} se dá na \textbf{inter-relação} da experimentação, decodificação e informação, na qual os arte-educadores \textbf{baseiam} -se na construção desse conhecimento e na premissa de \textbf{que} só \textbf{um} saber consciente e informado torna possível a aprendizagem significativa em estética-crítica-reflexiva em Arte.
A proposta da BNCC do Ensino Médio destaca os processos criativos como operadores de conhecimento estético, cultural, social e político, posto \textbf{que} Arte é \textbf{uma} área do conhecimento e não \textbf{uma} ferramenta e/ou instrumento \textbf{que} só se legitima se colocada a serviço de outras epistemologias e/ou saberes.
A Arte é propulsora de estudo, pesquisa e referências estéticas, poéticas, sociais, culturais e políticas, para criar novas relações entre \textbf{sujeitos} e seus modos de olhar para si e para o mundo ( BRASIL, 2018 ).
A aprendizagem da Arte deve ser \textbf{vista} para além de códigos e técnicas, e a produção artística deve ser observada como meio e não como fim.
Assim, o trabalho \textbf{didático-pedagógico} no ambiente escolar deve proporcionar ao estudante a protagonização vivenciada através do processo de articulação e criação.
Portanto, a prática docente tende a acentuar -se no planejamento coletivo e vai assegurar para o arte-educador e aos estudantes \textbf{uma} formação cultural na construção de sua prática pedagógica.
Importante ressaltar o \textbf{que} \textbf{diz} Iavelberg ( 2003, p. 09 ) acerca da inserção da Arte no contexto escolar quando afirma \textbf{que} “ a inserção não se dá apenas pela sua habilidade de desenvolver competências, habilidades e conhecimentos necessários a outros estudos, mas o seu valor intrínseco como construção humana, como patrimônio comum a ser apropriado por todos ”.
Nesse viés, a produção artística e cultural, com autonomia e características próprias de cada linguagem artística, atua na dimensão da potencialização da expressividade do estudante, considera o emocional individual ou coletivo do fenômeno do sentido da criação artística.
Vale \textbf{lembrar} \textbf{que} todo esse processo de articulação dentro desse campo de saber não acontece de forma espontânea, mas interdisciplinar, com \textbf{um} saber consciente para promover a compreensão e a fruição da Arte para todos.
Nessa perspectiva, a Arte pode ser considerada propulsora da interação \textbf{dialógica}, desde \textbf{que} seja articulada sempre como área de conhecimento autônoma, sendo respeitados os diversos princípios operativos de cada linguagem ( Artes Visuais, Dança, Teatro e Música ).
O campo da Arte dialoga com outras áreas do conhecimento, mas também é possível pensar em diálogos interdisciplinares entre as próprias linguagens artísticas, \textbf{uma} vez \textbf{que} cada \textbf{uma} tem elementos e objetos de estudo próprios, como se verifica ao abordar cada linguagem artística nesse documento.
Compreendemos \textbf{que} a \textbf{interdisciplinaridade} acontece quando há \textbf{um} aprendizado integrado e colaborativo não identificado com a polivalência.
Portanto, se faz necessário compreender \textbf{que} : > [... ] várias \textbf{disciplinas} se reúnem a partir de \textbf{um} mesmo objeto, porém é necessário criar -se \textbf{uma} situação-problema no sentido de Freire ( 1974 ), onde a ideia do projeto nasça da consciência comum, da fé dos investigadores no reconhecimento da complexidade do mesmo e na disponibilidade destes em redefinir o projeto a cada dúvida ou cada resposta encontrada ( FAZENDA, 2008, p. 22 ).
Desse modo, no desenvolvimento de projetos interdisciplinares, compreendemos \textbf{que} há a necessidade de \textbf{que} a Arte seja desenvolvida na escola nas suas diversas manifestações, e quando há \textbf{um} maior conjunto de habilidades \textbf{que} possam ser desenvolvidas pelas diferentes formações dos profissionais em Arte, diversas experiências e linguagens artísticas são realizadas pelos professores-arte-educadores, \textbf{que}, por meio do aprofundamento prático e \textbf{teórico}, possibilitam outras potencialidades sobre a \textbf{interdisciplinaridade} dos alunos no processo de \textbf{ensino-aprendizagem}, isto é, quando se expande o contato com demais formas de ensino da Arte, há \textbf{uma} formação integral e a aprendizagem significativa do \textbf{sujeito} do ensino médio.
A Arte na Educação como expressão individual ou coletiva e como cultura é \textbf{um} saber fundamental para a recognição cultural e o desenvolvimento pessoal, posto \textbf{que} na \textbf{divisão} dos saberes epistemológicos do ocidente não se pode categorizar a Arte como ferramenta.
Isto \textbf{posto}, a recognição cultural na Amazônia, histórias e imaginários estão repletos de saberes regionais \textbf{que} atestam \textbf{que} a Arte engendra conhecimento em \textbf{dialética} com outras áreas do conhecimento e, por isso, propicia aprendizado rico em diversidade cultural e sociabilidades \textbf{que} respeitem e problematizem as relações em sociedade.
Nessas histórias e imaginários, a Amazônia remonta ao inferno verde[61 ], à religiosidade, ao repouso, à encantaria, aos rios, às vicinais, aos campos, às florestas, às estradas, entre outras histórias \textbf{que} \textbf{permeiam} o sistema da Arte na Amazônia[62 ].
E o Ensino de Arte no Pará se legitima no Brasil, por meio das nuances de sentidos de opostos/complementares : real e imaginário de transrealidade[63 ].
Em Loureiro ( 1999 ), observa -se \textbf{que} a cultura provém da realidade de cada \textbf{um}, entrelaçada aos contextos culturais de memória e embebida em imaginários ( histórias orais ) em \textbf{que} os povos da floresta são distintos, mas com ampla visão de mundo.
Desse modo, o \textbf{autor} \textbf{diz} \textbf{que} : > A cultura é matéria em \textbf{que} o artista modula sua criação, \textbf{uma} vez \textbf{que} é através dessa ambiência \textbf{que} o \textbf{homem} \textbf{vive} e transvive a realidade.
E o \textbf{sumo} da cultura paraense-amazônica, creio, é \textbf{uma} forma de realismo mágico, \textbf{que} nos leva a entender a realidade por via de \textbf{uma} espécie de lógica mítica. [... ] A outra está no imaginário.
A cultura paraense – caboclo/ribeirinha – expressão do \textbf{que} chamo de Amazônia profunda – é fruto dessa coincidentia oppositorum entre real e imaginário.
O real nos coloca diante da necessidade prática de \textbf{viver}.
O imaginário nos garante o direito de sonhar ( LOUREIRO, 1999, p. 19, \textbf{grifos} do \textbf{autor} ).
As Artes Visuais, a Música, a Dança e o Teatro exercem \textbf{uma} precisão na originalidade, criatividade, técnica e sentido universal em suas especificidades.
Os artistas produzem para além de \textbf{um} lugar e são legitimados ao convívio nacional e internacional.
Portanto, vale \textbf{lembrar} \textbf{que} todo artista cria apoiado em \textbf{uma} herança cultural local e universal ; no entanto, ao propor o diálogo entre a criação e a herança cultural na Amazônia, este documento curricular apresenta \textbf{uma} relação \textbf{dialética} da criação individual com o conjunto de valores ( materiais e espirituais ) universais, \textbf{que} ressalta os conhecimentos \textbf{antropológicos} e histórico-social dos \textbf{sujeitos} amazônidas[64 ] acumulados ao longo dos tempos.
Dessa forma, a autonomia criativa e \textbf{expressiva} se apresenta como outra forma de conhecimento, apreensão e compreensão do mundo \textbf{que} se coloca para além da \textbf{simples} racionalidade técnico-científica sem, no entanto, deixar de contribuir com habilidades e competências específicas das outras ciências como, por exemplo : a Música e as Artes Visuais com a Matemática ou a Dança e o Teatro com a Anatomia.
Essas \textbf{inter-relações} habilitam os estudantes para \textbf{um} saber consciente e informado, com o objetivo de tornar possível a aprendizagem em Arte.
Barbosa ( 2002 ) destaca \textbf{que} : > Não mais se pretende desenvolver \textbf{uma} vaga sensibilidade nos alunos por meio da Arte, mas também se aspira influir positivamente no desenvolvimento cultural dos estudantes pelo ensino/aprendizagem da Arte.
Não podemos entender a Cultura de \textbf{um} país sem conhecer sua Arte.
A arte como \textbf{uma} linguagem aguçadora dos sentidos transmite significados \textbf{que} não podem ser transmitidos por nenhum outro tipo de linguagem, tais como a discursiva e a científica ( BARBOSA, 2002, p. 17-18 ).
Portanto, os aspectos corporais, gestuais, teatrais, visuais, espaciais e sonoros consideram particularidades e instituem aos estudantes conhecer as especificidades dos saberes e práticas do ensino da Arte, \textbf{que}, por meio de suas linguagens ( Artes Visuais, Teatro, Música e Dança ), têm \textbf{um} professor para cada linguagem.
Este trabalha estéticas de diversas manifestações artísticas e culturais, assim como as associações entre elas, de maneira interdisciplinar.
A \textbf{inter-relação} dessas linguagens propicia processos criativos para a construção das pesquisas e estudos de referências estéticas, poéticas, sociais, culturais e políticas, bem como a criação de novas relações entre os \textbf{sujeitos} e seus modos de olhar para si e para o mundo, no sentido de desconstruir para reconstruir, selecionar, reelaborar, a partir do \textbf{conhecido} e modificá -los de acordo com o contexto e a necessidade, desenvolvidos pelo ler, fazer e contextualizar Arte, essenciais para a vida cotidiana comprometida com a diversidade cultural na Amazônia e no mundo.
O engajamento com a diversidade cultural é enfatizado pela Arte-Educação[65 ], e essa participação ativa referente às questões políticas e sociais foi consolidada a partir do movimento de Arte Moderna, com a proposta de diversidade de códigos em função de etnias, \textbf{raças}, gênero, classe social, entre outros, e não mais, especificamente, nos códigos europeus e norte-americanos.
Esse movimento é considerado a grande ruptura no modo de conceber a Arte e o seu ensino, \textbf{que} tradicionalmente era centralizada no ensino da técnica.
Barbosa ( 1975, p. 44 ) estabelece \textbf{que} “ na realidade, nossa primeira grande renovação \textbf{metodológica} no campo da Arte-Educação se deve ao movimento de Arte Moderna de 1922”[66 ].
A relevância da Arte como conhecimento surge em contraposição às teses Liberais, Positivistas e Modernistas e a \textbf{preconiza} como \textbf{uma} construção social, histórica e cultural aliada à dimensão da \textbf{cognição}.
Esse processo construtor da Arte ligada à \textbf{cognição} é relevante porque alia -se à concepção da razão e rege a Arte como \textbf{um} campo de conhecimento com suas estruturas e bases \textbf{que} se fundamentam à \textbf{cognição}.
Rizzi ( 2002, p. 64-65 ) propõe \textbf{uma} relação ao propósito da Arte como Conhecimento frente à Arte na Educação e a Arte em si, na medida em \textbf{que} “ acredita ser a Arte importante por si mesma e não por ser instrumento para fins de outra natureza.
Por ser \textbf{uma} experiência \textbf{que} permite a integração da experiência singular e isolada de cada ser humano com a experiência da \textbf{humanidade} ”.
Assim, essa concepção \textbf{que} congrega a experiência \textbf{baseada} na idiossincrasia aliada à experiência da \textbf{humanidade} requer \textbf{uma} integração fundamentada no interculturalismo, na \textbf{interdisciplinaridade} e na aprendizagem dos conhecimentos artísticos, a partir da \textbf{inter-relações} entre História da Arte/Contextualização, Leitura da Imagem – Apreciação Estética e Fazer Artístico.
Embora alguns \textbf{autores} utilizem termos como multiculturalismo e pluriculturalidade para definir diversidade cultural, este documento sugere a utilização do termo “ Intercultural ” de Fleuri ( 2003 ) : > A intercultura refere -se a \textbf{um} complexo campo de debate entre as variadas concepções e propostas \textbf{que} enfrentam a questão da relação entre processos identitários socioculturais diferentes, focalizando especificamente a possibilidade de respeitar as diferenças e de integrá -las em \textbf{uma} unidade \textbf{que} não as anule.
A intercultura vem se \textbf{configurando} como \textbf{uma} nova perspectiva epistemológica, ao mesmo tempo é \textbf{um} objeto de estudo interdisciplinar e transversal, no sentido de tematizar e teorizar a complexidade ( para além da pluralidade ou da diversidade ) e a ambivalência ou o hibridismo ( para além da reciprocidade ou da evolução linear ) dos processos de elaboração de significados nas relações intergrupais e intersubjetivas, constituídas de campos identitários em \textbf{termos} de etnias, de gerações, de gênero, de ação social ( FLEURI, 2003, p. 17 ).
O conceito de interculturalidade contribui para o aprofundamento de debates em função de sua complexidade.
Entretanto, a proposição acerca da interação entre as diferentes culturas alia -se aos objetivos da Arte-Educação \textbf{comprometida} com o desenvolvimento cultural, com o desenvolvimento da crítica, com o conhecimento acerca da cultura do país e, fundamentalmente, da cultura \textbf{que} nos é peculiar.
A Arte-Educação alicerçada na Comunidade[67 ] é \textbf{uma} proposição \textbf{contemporânea} \textbf{que} tem indicativos positivos em projetos de educação com o ponto de atenção na reconstrução social, e dá lugar ao em-comum, com o pensamento na essência e existência.
Esse pensamento, entre linguagem e mundo, não isola a cultura local, mas a discute em relação com outras culturas.
Impulsionar a \textbf{inter-relação} de culturas e saberes para promover aos estudantes o acesso às diversas manifestações culturais \textbf{populares} presentes na comunidade é \textbf{desejável}.
A diversidade das populações tradicionais ( indígenas, \textbf{ribeirinhos}, quilombolas, caiçaras, citadinos e outros grupos ), a ecologia, a sustentabilidade, o progresso, o conhecimento, a justiça social, a geopolítica e tantos outros processos nos âmbitos socioculturais tornam a Amazônia intercultural.
A \textbf{influência} dos povos indígenas na gastronomia regional visionada no contexto nacional e reconhecida internacionalmente, assim como a sonoridade do Pará ( carimbó, brega, toada do boi, entre outros ), expressam a identidade cultural, artística, social, ambiental e histórica da região amazônica.
Paralelamente, nesse diálogo participante, ressaltam -se os centros culturais como os museus, os teatros, as \textbf{exposições} de arte, os recitais de música, as apresentações de dança em espaços públicos, os festivais, as experiências interdisciplinares e interculturais, colocando \textbf{lado} a \textbf{lado} o código erudito e o \textbf{popular} como, por exemplo, os especiais saberes dos mestres ceramistas, as mulheres erveiras[68 ], entre outros \textbf{que} acontecem na capital Belém e em todas as regiões da Amazônia.
Vale \textbf{lembrar} \textbf{que} as Leis 10.639/03[69 ] e 11.645/2008[70 ] impulsionam \textbf{um} olhar para além da produção artística ocidental, de matriz europeia, e integram a obrigatoriedade das culturas indígenas, africanas e afro-brasileiras no Ensino Fundamental e Médio.
Esse ponto é relevante, pois \textbf{qualifica} o ensino-aprendizagem acerca do reconhecimento e valorização das diversas culturas \textbf{que} constituem o Brasil.
Por esse viés, se alcança \textbf{uma} ligação \textbf{imediata} com o universo social do estudante.
As diferenças são valorosas para promover a constituição do \textbf{sujeito}, levando -se em conta \textbf{uma} cultura voltada para o envolvimento da comunidade escolar relacionada às manifestações artísticas e culturais paraenses, aliadas às artes consagradas mundialmente.
Essa ampliação dar -se-á pelo conceito de interculturalidade.
Sobre esse apontamento, Barbosa ( 1995, p. 11 ) ressalta \textbf{singularidades} importantes para a construção de \textbf{um} currículo crítico ao afirmar \textbf{que} : > No \textbf{que} \textbf{diz} respeito à cultura local, pode -se constatar \textbf{que} apenas o nível erudito desta cultura é admitido na escola.
As culturas das classes sociais baixas continuam a ser ignoradas pelas instituições educacionais, mesmo pelos \textbf{que} estão envolvidos na educação destas classes.
Nós aprendemos com Paulo Freire ( 1987 ) a rejeitar a segregação cultural na educação.
As décadas de luta para salvar os oprimidos da ignorância sobre eles próprios nos ensinaram \textbf{que} \textbf{uma} educação libertária terá sucesso só quando os participantes no processo educacional forem capazes de identificar seu ego cultural e se orgulharem dele.
Isto não significa a defesa de guetos culturais, nem excluir a cultura erudita das classes baixas.
Todas as classes sociais têm o direito de acesso aos códigos da cultura erudita porque esses são os códigos dominantes - os códigos do poder.
É necessário conhecê -los, ser versados neles, mas tais códigos continuarão a ser \textbf{um} conhecimento exterior a não ser \textbf{que} o indivíduo tenha dominado as referências culturais da sua própria classe social, a porta de entrada para a \textbf{assimilação} do"outro".
A mobilidade social \textbf{depende} da \textbf{inter-relação} entre os códigos culturais das diferentes classes sociais.
Essa mediação entre cultura local e global / erudita e \textbf{popular} dar -se-á pelos debates \textbf{que} \textbf{permeiam} o conceito de interculturalidade, já apresentado anteriormente, \textbf{que} amplia o exercício da crítica, da apreciação e da fruição num sentido contextualizado e corresponde à epistemologia da arte e amplia aos modos como se aprende Arte.
Pois, como \textbf{lembra} Barbosa ( 2002, p. 19 ), para alcançar esse objetivo com qualidade, é “ necessário \textbf{que} a escola forneça \textbf{um} conhecimento sobre a cultura local, a cultura de vários grupos \textbf{que} caracterizam a nação e a cultura de outras nações ”.
Nesse sentido, a arte \textbf{contemporânea} \textbf{atravessou} significativamente o momento em \textbf{que} a cultura ganhou maior expressão econômica e social no mundo.
Isso \textbf{posto}, Ramos ( 2006 ) ressalta essa importância, ao \textbf{dizer} \textbf{que} : > A Arte \textbf{contemporânea} pôde ser entendida como sinal e instrumento ativo de aprendizagem ; por estimular a troca simbólica entre grupos e culturas diferentes, por estimular a \textbf{percepção} para novos conteúdos e nova leitura de vida social.
Na verdade, tudo \textbf{que} está fora do sistema da arte pode ser incorporado por ela, passando a fazer parte de seu jogo.
Dessa maneira, as galerias de arte, os museus e centros culturais são importantes formadores de capital simbólicos, aliados – pelos interesses mais nobres e mais vis – a agentes e meios de comunicação capazes de transformar significativamente as sensibilidades e conteúdos do público ao seu alcance ( RAMOS, 2006, p. 22 ).
Assim, a linguagem artística articula -se na perspectiva da \textbf{interdisciplinaridade} com os outros campos de saberes e práticas de Ensino e outros campos de Saberes e práticas Eletivos, assim como entre as linguagens da Arte capazes de promover a liberdade de escolhas acerca de práticas pedagógicas como dimensão didática do processo ensino-aprendizagem.
Nesse sentido, as \textbf{abordagens} \textbf{metodológicas} específicas do campo das Artes Visuais, Teatro, Música e da Dança são específicas e orientam as \textbf{abordagens} para os docentes de tais campos conforme se verifica nas Linguagens[71 ].
As Artes Visuais são \textbf{um} conjunto de manifestações \textbf{que} se iniciam quando o \textbf{homem} produziu objetos utilitários e fez pinturas ou gravuras nas rochas ; naquele momento, a \textbf{humanidade} busca “ [... ] tornar visível aquilo \textbf{que} aponta para a existência de pensamentos invisíveis : ou seja, símbolos. ” ( BELL, 2008, p. 10 ).
Nesse contexto, o ensino/aprendizagem das Artes Visuais no Ensino Médio se mostra fundamental para a formação de \textbf{um} estudante/cidadão consciente e capaz de se tornar \textbf{um} leitor crítico sobre as imagens \textbf{que} constituem o mundo ao seu redor.
O \textbf{sujeito} do ensino médio no contexto proposto neste documento, por meio do estudo das Artes Visuais, visa à compreensão da multiplicidade \textbf{que} compõe as várias expressões visuais ao longo de toda a existência humana.
É importante ressaltar a \textbf{abordagem} Triangular através do vértice Ler, Fazer e Contextualizar, \textbf{que} surge no campo das Artes Visuais e pode ser experimentada em outras linguagens artísticas, com \textbf{encaminhamentos} \textbf{metodológicos} diferentes devido às especificidades \textbf{que} apresentam, conforme se verifica no decorrer do texto, com as linguagens de Teatro, Dança e Música.
Tais \textbf{encaminhamentos} têm o objetivo de \textbf{focar} na proposição escolhida pelo professor nas suas aulas práticas, sem vínculo \textbf{teórico} padronizado, a fim de não engessar o processo.
Para Barbosa ( 2010, p. 10 ), “ refere -se à \textbf{uma} \textbf{abordagem} eclética.
Requer transformações enfatizando o contexto ”.
Barbosa ( 2002, p. 18 ) enfatiza ainda \textbf{que} “ por meio da Arte é possível desenvolver a \textbf{percepção} e a imaginação, apreender a realidade do meio ambiente, desenvolver a capacidade crítica, permitindo ao indivíduo analisar a realidade percebida e desenvolver a criatividade de maneira a mudar a realidade \textbf{que} foi analisada ”.
Nesse propósito, a Arte na Educação atribui \textbf{um} \textbf{raciocínio} crítico como expressão pessoal e como cultura, para desenvolver a \textbf{percepção} e a imaginação diante da complexidade do mundo, promovendo \textbf{um} diálogo intercultural ( pluralidade ) \textbf{que} estabelece na sala de aula o exercício de sua cidadania nas dimensões estéticas, política e social, posto \textbf{que} o exercício da Arte não se encontra fora da esfera da cidadania, mas sim é parte e está \textbf{dialeticamente} relacionado a ela.
Nesse contexto, o ensino/aprendizagem das Artes Visuais no Ensino Médio se mostra fundamental para a formação de \textbf{um} estudante/cidadão consciente e capaz de se tornar \textbf{um} leitor crítico sobre as imagens \textbf{que} constituem o mundo ao seu redor.
Diante desta proposta, entendemos as conexões das áreas de conhecimento apontadas com os objetos de conhecimento da área de Linguagens e Suas Tecnologias, visando à superação dos modelos espontaneístas e tradicionalistas de Ensino da Arte, e, assim, possibilitar a autonomia dos professores na construção de suas propostas das suas metodologias pedagógicas. [ 61 ] : O termo “ inferno verde ” foi corriqueiramente empregado para designar a região amazônica, como alusão às difíceis condições de vida na floresta durante a colonização.
Esse termo foi aprofundado por Euclides da Cunha com sua obra \textbf{Um} paraíso perdido.
Ensaios amazônicos, em \textbf{que} lança \textbf{um} olhar especial sobre a região, ao chefiar a Comissão Brasileira de Limites com o Peru, em 1905.
Cunha também irá prefaciar a obra Inferno Verde - scenas e scenario do Amazonas, de Alberto Rangel, 1907 ( MANESCHY, 2010, p. 11 ). [ 62 ] : Conhecer a Arte na região Amazônica, acompanhar a apreensão \textbf{que} os artistas têm de seus territórios e as modulações de seus discursos é \textbf{tarefa} para quem busca contribuir à construção de \textbf{um} entendimento do país \textbf{que} é diverso e multifacetado. [... ] Artistas vêm organizando complexas perspectivas de entendimento acerca de seus territórios - físicos, culturais, políticos, estéticos - e \textbf{que} têm elevado seus olhares para outros lugares, transcendendo regionalismo ( MANESCHY, 2010, p. 12 ). [ 63 ] : Loureiro, ao usar o termo “ transrealidade ”, refere \textbf{que}, no Pará, como em toda a Amazônia, há \textbf{um} conluio de transracionalidade, de transgressão da lógica dualística, do maravilhamento ardente ( incompleto ).
Somos surreais como os gregos eram míticos ( LOUREIRO, 1999, p. 19 ). [ 64 ] : Os \textbf{sujeitos} amazônidas são frutos da confluência de \textbf{sujeitos} sociais distintos – ameríndios da várzea e/ou terra firme, negros, nordestinos e europeus de diversas nacionalidades ( portugueses, espanhóis, holandeses, franceses, etc. ) – \textbf{que} inauguram novas e singulares formas de organização social nos trópicos amazônicos.
Diferenciada em suas matrizes geracionais, marcada por dinamismos e sincretismos singulares, a formação social amazônica foi fundamentada \textbf{historicamente} em tipos variados de escravismo e servidão.
Assim, falar dos povos da Amazônia requer \textbf{um} (re)conhecimento da grande diversidade ambiental e social da região ; noutras palavras, é \textbf{preciso} tomar como ponto de partida o desenvolvimento histórico da região ( FRAXE ; WITKOSKI ; MIGUEZ, 2009, p. 1 ). [ 65 ] : O termo Arte-educação ( tradução do inglês art-education ) surgiu como \textbf{uma} tentativa de conectar e integrar os conceitos de arte e educação, resgatando suas relações mais significativas e reforçando a ideia de continuidade entre seus \textbf{fundamentos}, contribuindo para transformar sentimentos e ideias.
Para Barbosa ( 1991, p. 07 ), a “ Arte-educação é epistemologia da arte e, portanto, é a investigação dos modos como se aprende arte na escola [... ], na universidade e na intimidade dos ateliers ”. [ 66 ] : É inegável a importância do movimento de Arte Moderna no Brasil, \textbf{que} teve como ponto de atenção as concepções de Ensino da Arte como expressão e também como atividade, porém a função ideológica, com o propósito de colocar no centro o caráter humanista no currículo, e a isenção de qualquer conteúdo de Arte no “ fazer artístico nas escolas ” trouxeram \textbf{um} esvaziamento ao ensino de Arte e, mais ainda : a colocou em \textbf{um} lugar de secundariedade na educação escolar. [ 67 ] : Para Agamben ( 2007 ), pensar a comunidade é antes de mais nada pensar \textbf{que} não há apenas a linguagem encerrada em si, produtora de enunciados auto-reflexivos e autotélicos, mas as coisas singulares, \textbf{que} nela se \textbf{perdem} através dessa perda no comum, no universal, ganham \textbf{uma} presença, \textbf{que} é produtora de interrupções, de vazios, pelos quais o comum \textbf{que} se faz pela linguagem nela se desfaz para dar lugar ao em-comum.
A comunidade \textbf{que} vem dá lugar no seu acontecer ; não se confunde com quaisquer comunidades reais, \textbf{atuais}, identificáveis.
Pelo contrário, é no desfazer destas \textbf{que} se afirma ; é \textbf{enquanto} apropriação/expropriação \textbf{que} a linguagem se apresenta condição do humano, e se anuncia \textbf{uma} época, a nossa, “ \textbf{que} não quer ser já é \textbf{uma} época histórica ” ( SEDLMAYER ; GUIMARÃES ; OTTE, 2007, p. 72 ). [ 68 ] : Chamadas erveiras ou cheirosinhas das feiras públicas de Belém do Pará nos mostram \textbf{que} esta é \textbf{uma} atividade eivada de mudanças e permanências, sendo elas as mais diversas, desde a mudança temporal até a mudança geográfica.
Isso, devido aos vários lugares em \textbf{que} essas personagens têm atuado ao longo do tempo até chegarem ao mercado do Ver-o-Peso, lugar onde \textbf{hoje} elas possuem \textbf{um} grau de importância \textbf{que} vai além da \textbf{simples} comercialização das ervas, chegando ao nível simbólico, religioso e turístico ( SILVA, 2017, p. 250 ). [ 69 ] : A Lei 10.639/03, alterada pela Lei 11.645/08, \textbf{que} torna obrigatório o ensino da história e cultura afro-brasileira e africana em todas as escolas, públicas e particulares, do ensino fundamental até o ensino médio. [ 70 ] : A Lei 11.645/2008 altera a Lei nº 9.394, de 20 de dezembro de 1996, modificada pela Lei nº 10.639, de 9 de janeiro de 2003, \textbf{que} estabelece as diretrizes e bases da educação nacional, para incluir no currículo oficial da rede de ensino a obrigatoriedade da temática “ História e Cultura Afro-Brasileira e Indígena ”. [ 71 ] : As contribuições específicas de Artes Visuais, Teatro, Música e Dança ( itens 1.1, 1.2, 1.3 e 1.4 ) foram extraídas das considerações elaboradas pelo Coletivo Pró-Associação de Arte Educadores do Estado do Pará - AAEPA.
Essas considerações foram feitas na consulta pública e estão acessíveis ao público no site da Federação de Arte Educadores do Brasil - FAEB.

% matched lemmas: abordagem, antropológico, assimilação, atravessar, atual, autor, basear, cognição, comprometido, configurar, conhecido, contemporâneo, depender, desejável, dialeticamente, dialético, dialógico, didático-pedagógico, disciplina, divisão, dizer, encaminhamento, enquanto, ensino-aprendizagem, exposição, expressivo, focar, fronteira, fundamento, grifo, historicamente, hoje, homem, humanidade, imediato, influência, inter-relação, interdisciplinaridade, lado, lembrar, metodológico, operar, percepção, perder, permear, popular, postar, preciso, preconizar, qualificar, que, raciocínio, raça, ribeirinho, simples, singularidade, sujeito, sumo, tarefa, termos, teórico, um, uma, ver, viver
\end{textsample}
